\documentclass[UTF8,a4paper,11pt,fontset=fandol]{ctexart}
\usepackage[margin=23mm]{geometry}
\usepackage{amsmath,amssymb,booktabs,longtable,graphicx,float,fancyhdr,xcolor,needspace}
\usepackage{xurl}
\usepackage[unicode,colorlinks=true,linkcolor=black,urlcolor=blue]{hyperref}
\definecolor{accent}{HTML}{155574}
\pagestyle{fancy}
\fancyhf{}
\fancyhead[L]{\small 数学建模训练论文}
\fancyhead[R]{\small 待人工审阅}
\fancyfoot[C]{\thepage}
\setlength{\headheight}{15pt}
\setlength{\parskip}{0.35em}
\setlength{\emergencystretch}{3em}
\title{\bfseries 自定义适配器示例：一维二次损失最小化}
\author{}
\date{}
\begin{document}
\maketitle
\vspace{-2.4em}
\begin{center}\small\color{accent}训练草稿：计算已核验，建模解释、引用及赛事格式仍需队员审阅\end{center}
\begin{abstract}
针对等权观测的中心估计问题，构建平方损失最小化模型，得到中心估计 5.00，平方损失为 20.00。通过梯度条件及邻域扰动检查验证结果。本例用于演示自定义适配器与统一论文构建链路。
\end{abstract}
\noindent\textbf{关键词：}自定义适配器；最小二乘；独立核验

\section{问题分析与数据来源}
人工构造的独立适配器示例数据，仅用于验证模型扩展接口。
\begin{description}

\item[Q1] 估计最小化平方损失的中心值，并独立核验一阶最优条件。

\end{description}

\section{模型假设与符号}
\begin{enumerate}

\item 各观测具有相同权重。

\end{enumerate}
\begin{longtable}{p{.16\textwidth}p{.46\textwidth}p{.20\textwidth}}
\toprule 符号 & 含义 & 单位 \\ \midrule

c & 中心值 & 单位 \\

\bottomrule
\end{longtable}

\section{模型建立与求解方法}
等权平方损失是严格凸函数，样本均值给出唯一极小值。

\begin{equation}
\min_c L(c)=\sum_i(y_i-c)^2
\end{equation}

\noindent\textbf{基线方案：}任意固定中心。


\section{模型求解与验证}
平方损失的一阶条件表明，最优中心为样本均值。本次结果为 5.00，对应损失为 20.00。核验代码重新计算梯度并检查两侧扰动的损失，避免只依赖求解函数的返回状态。


\section{计算结果与图表}
所有数值均读取同一次实验记录；表中的不可用项保留其原始状态。

\Needspace{7\baselineskip}
\subsection*{中心估计结果}
{\small
\begin{longtable}{ll}
\toprule
中心值 & 观测数 \\ 
\midrule\endhead

5 & 4 \\ 

\bottomrule
\end{longtable}}



\section{检验、局限与改进方向}
本次自动核验覆盖了已编码的数值条件。自动核验通过不等价于题意解释已经充分，队员仍需对照原题审阅假设与结论。
\begin{itemize}

\item 平方损失对离群值敏感。

\end{itemize}

\section*{参考来源}

本例未使用外部经验事实；模型由所列数学关系直接推导。实际赛题使用外部资料时，应在此补充并核验来源。

\section*{复现与 AI 使用说明}
数据来源、程序版本和输入摘要见工作区中的 provenance.json；运行命令见复现说明。
本案例使用已声明的训练数据，结论仅用于演示建模过程。
实际 AI 工具使用、人工修改及核验记录保存在 ai\_usage.jsonl 中，由队员据实补充后按赛事要求整理。
\end{document}